\documentclass[11pt]{article}

\usepackage[a4paper,margin=1in]{geometry}
\usepackage{amsmath,amssymb,amsthm,mathtools}
\usepackage[hidelinks]{hyperref}

\newtheorem{theorem}{Theorem}
\newtheorem{proposition}[theorem]{Proposition}
\newtheorem{lemma}[theorem]{Lemma}
\newtheorem{corollary}[theorem]{Corollary}
\theoremstyle{definition}
\newtheorem{remark}[theorem]{Remark}

\newcommand{\R}{\mathbb{R}}
\newcommand{\eps}{\varepsilon}
\newcommand{\wt}{\widetilde}
\newcommand{\Om}{\Omega}
\newcommand{\dx}{\,dx}

\title{The Quantitative Single-Set Selection Principle\\for Sharp Faber--Krahn Stability (Route A)}
\author{}
\date{}

\begin{document}
\maketitle

\begin{abstract}
This note is the building block \textsc{SelectionPrinciple} of the Route~A
program for sharp Faber--Krahn stability with exponent~$1/2$,
\[
\Bigl(\frac{|\Om|}{|B_1|}\Bigr)^{2/N}
\bigl(\lambda_1(\Om)-\lambda_1(B^*)\bigr)
\ge c_{\rm FK}(N)\,\mathcal A(\Om)^2.
\]
We isolate one quantitative step of the Brasco--De~Philippis--Velichkov
proof~\cite{BDV} (BDV): given one near-counterexample $\Om_0$ with small
Saint--Venant quotient, we produce a \emph{single} selected quasi-open set
$\wt\Om$ of unit volume whose energy-to-asymmetry quotient is controlled by
that of $\Om_0$, on which the BDV free-boundary regularity package applies,
and which carries an explicit volume-rescaled Bernoulli law on
$\partial\wt\Om$.

The proof contains no compactness or contradiction step. All BDV results
are quoted as quantitative black boxes whose constants depend only on $N$
and $R$ (and are extractable from BDV, but not made numerically explicit
here).
\end{abstract}

\section{Setting and Normalisation}

Fix $N\ge2$ and $R\ge2$. We use the BDV torsion convention: for a
quasi-open set $\Om\subset B_R\subset\R^N$,
\[
E(\Om)
:=\min_{u\in H^1_0(\Om)}\!
\Bigl\{\tfrac12\!\int_\Om|\nabla u|^2-\!\int_\Om u\Bigr\}
=-\tfrac12\!\int_\Om u_\Om,
\]
where $u_\Om$ is the torsion function. The Saint--Venant deficit is
\[
\delta(\Om):=E(\Om)-E(B_1)\ge0
\qquad\text{whenever }|\Om|=|B_1|=\omega_N.
\]

\subsection*{Asymmetries.}
Let $\mathcal A(\Om)$ be the Fraenkel asymmetry. Following BDV we work with
the smoother asymmetry
\[
\alpha(\Om):=\int_{\Om\Delta B_1(x_\Om)}\bigl|1-|x-x_\Om|\bigr|\dx,
\qquad x_\Om:=\frac1{|\Om|}\!\int_\Om x\dx.
\]
BDV Lemma~4.2 gives, for $\Om\subset B_R$,
\begin{align}
|\Om\Delta B_1(x_\Om)|^2 &\le C_1(N)\,\alpha(\Om), \label{eq:42i}\\
|\alpha(\Om_1)-\alpha(\Om_2)| &\le C_2(R)\,|\Om_1\Delta\Om_2|. \label{eq:42ii}
\end{align}
Throughout we set, with a small abuse of notation matching downstream
blocks,
\[
A(\Om):=\alpha(\Om),\qquad
Q_\alpha(\Om):=\frac{E(\Om)-E(B_1)}{\alpha(\Om)}.
\]

\subsection*{Penalised functional.}
Let $f_{\hat\eta}$ be the BDV one-sided volume penalty from Lemma~4.5,
\[
f_{\hat\eta}(s)=
\begin{cases}
\hat\eta(s-\omega_N),&s\le\omega_N,\\
\hat\eta^{-1}(s-\omega_N),&s\ge\omega_N,
\end{cases}
\]
with $\hat\eta=\hat\eta(N,R)\in(0,1)$. Set
$F_{\hat\eta}(U):=E(U)+f_{\hat\eta}(|U|)$. BDV Lemma~4.5 gives that $B_1$
minimises $F_{\hat\eta}$ over quasi-open $U\subset B_R$, and the linear
bilateral bound
\begin{equation}\label{eq:Lem45}
F_{\hat\eta}(U)-F_{\hat\eta}(B_1)\ge C_4(N,R)^{-1}\bigl||U|-|B_1|\bigr|
\end{equation}
holds.

Given an input $\Om_0\subset B_R$ with $|\Om_0|=|B_1|$, write
$\eps:=\alpha(\Om_0)$, $\delta:=E(\Om_0)-E(B_1)$, $q:=\delta/\eps$, and
choose $\tau:=q^{1/4}$. Define
\begin{equation}\label{eq:Jdef}
J(U):=F_{\hat\eta}(U)
+\sqrt{\eps^2+\tau^2(\alpha(U)-\eps)^2}.
\end{equation}
This is BDV's penalised functional (4.10) with the BDV sequence parameter
$\sigma$ replaced by a single-set parameter $\tau$ tied to the input
deficit.

\subsection*{Thresholds.}
All depending only on $N,R$:
\begin{itemize}
\item $\tau_{\rm ex}(N,R)\le\hat\eta/(2C_2)$, ensuring existence of a
minimiser of $J$ and the uniform perimeter bound of BDV Lemma~4.6;
\item $\sigma_2(N,R)$ from BDV Lemma~4.9 (density estimates);
\item $\sigma_3(N,R)$ from BDV Lemma~4.16 (smooth Bernoulli coefficient);
\item $\tau_{\rm reg}(N,R):=\min\{\tau_{\rm ex},\sigma_2,\sigma_3\}$.
\end{itemize}
We assume
\begin{equation}\label{eq:Sml}
q\le q_{\rm sel}(N,R):=\bigl(\tau_{\rm reg}(N,R)\bigr)^{4},\qquad
\eps\le\eps_{\rm sel}(N,R).
\end{equation}
The second smallness keeps the dilated set inside a fixed enlarged ambient
ball. Under~\eqref{eq:Sml}, $\tau=q^{1/4}\le\tau_{\rm reg}$.

\section{Existence by the BDV Direct Method}\label{sec:existence}

\begin{lemma}[Existence + uniform perimeter; BDV Lemma~4.6, single-set form]
\label{lem:ex}
For every $\tau\le\tau_{\rm ex}(N,R)$, $J$ admits a quasi-open minimiser
$U_*\subset B_R$, $|U_*|>0$, with
\[
P(U_*)\le P_{\rm sel}(N,R).
\]
\end{lemma}

\begin{proof}[Why this is not a hidden compactness step.]
The BDV existence proof for (4.10) applies the direct method to the
truncated lower-semicontinuous part $U\mapsto F_{\hat\eta}(U)$, using
\[
\sqrt{\eps^2+\tau^2(\alpha(U)-\eps)^2}\ge\tau\,|\alpha(U)-\eps|
\]
together with~\eqref{eq:42ii} to bound $\alpha(U)$ along a minimising
sequence and force $|U|\le M(N,R)$. The closure step uses
$\gamma$-convergence of capacitary measures inside $B_R$ as in BDV~\S2,
whose constants depend only on $N,R$. Importantly, the limit set is
admitted as a minimiser without any selection from a sequence of
\emph{counterexamples}: the input here is a single $\Om_0$ and existence is
for the functional $J$ at fixed parameters $(\eps,\tau)$.

The uniform perimeter bound follows from BDV (4.13)--(4.14): testing
$J(U_*)\le J(\Om_0)\le F_{\hat\eta}(\Om_0)+\eps$ and using the
isoperimetric inequality with~\eqref{eq:Lem45} produces a bound depending
only on $N,R$.
\end{proof}

We fix henceforth a minimiser $U_*$ from Lemma~\ref{lem:ex}.

\section{The Three Quantitative Lemmas}

\subsection*{Lemma 1. Energy and asymmetry preservation before normalisation.}

\begin{lemma}\label{lem:1}
The minimiser $U_*$ satisfies
\begin{equation}\label{eq:L1a}
0\le F_{\hat\eta}(U_*)-F_{\hat\eta}(B_1)\le\delta,
\end{equation}
\begin{equation}\label{eq:L1b}
|\alpha(U_*)-\eps|\le\frac{\sqrt{2\eps\delta+\delta^2}}{\tau}.
\end{equation}
In particular, with $\tau=q^{1/4}$ and $q\le1$,
$|\alpha(U_*)-\eps|\le 2\eps\,q^{1/4}.$
\end{lemma}

\begin{proof}
By minimality $J(U_*)\le J(\Om_0)=F_{\hat\eta}(\Om_0)+\eps$. Since
$|\Om_0|=|B_1|$, $F_{\hat\eta}(\Om_0)=E(\Om_0)=F_{\hat\eta}(B_1)+\delta$.
As $B_1$ minimises $F_{\hat\eta}$, $F_{\hat\eta}(U_*)\ge F_{\hat\eta}(B_1)$.
Hence
\[
F_{\hat\eta}(U_*)+\sqrt{\eps^2+\tau^2(\alpha(U_*)-\eps)^2}
\le F_{\hat\eta}(B_1)+\delta+\eps,
\]
giving \eqref{eq:L1a} and
$\sqrt{\eps^2+\tau^2(\alpha(U_*)-\eps)^2}\le\eps+\delta$. Squaring yields
\eqref{eq:L1b}.
\end{proof}

\subsection*{Lemma 2. Volume error.}

\begin{lemma}\label{lem:2}
There is $C_5=C_5(N,R)$ such that
\[
\bigl||U_*|-|B_1|\bigr|\le C_5\,\delta,\qquad
|r_*-1|\le C_5\,\delta,\quad r_*:=(|U_*|/|B_1|)^{1/N}.
\]
\end{lemma}

\begin{proof}
By Schwarz symmetrisation, $F_{\hat\eta}(B_{r_*})\le F_{\hat\eta}(U_*)$
where $B_{r_*}$ is a ball of volume $|U_*|$. Combining with~\eqref{eq:Lem45}
and~\eqref{eq:L1a},
\[
C_4^{-1}\bigl||U_*|-|B_1|\bigr|\le F_{\hat\eta}(B_{r_*})-F_{\hat\eta}(B_1)
\le F_{\hat\eta}(U_*)-F_{\hat\eta}(B_1)\le\delta.
\]
Hence $||U_*|-|B_1||\le C_4\delta$, and bilipschitz dependence of $r_*$
on $|U_*|$ near $1$ yields $|r_*-1|\le C_5\delta$.
\end{proof}

\subsection*{Lemma 3 (fully expanded). Quotient preservation after volume normalisation.}

Define $\wt\Om:=r_*^{-1}U_*$, $r_*=(|U_*|/|B_1|)^{1/N}$.

\begin{lemma}\label{lem:3}
Under~\eqref{eq:Sml}, with $\tau=q^{1/4}$, after translation by $-x_{U_*}$
if needed:
\begin{enumerate}
\item[\rm(a)] $|\wt\Om|=|B_1|=\omega_N$.
\item[\rm(b)] $|\alpha(\wt\Om)-\eps|\le\rho\,\eps$ with
$\rho=2q^{1/4}+C_\alpha(N,R)\,q\le1/2$.
\item[\rm(c)] $E(\wt\Om)-E(B_1)\le C_E(N,R)\,\delta$.
\item[\rm(d)] $\displaystyle Q_\alpha(\wt\Om)\le\frac{C_E(N,R)}{1-\rho}\,Q_\alpha(\Om_0).$
\end{enumerate}
\end{lemma}

The proof rests on three quantitative inputs.

\paragraph{Step 1: Lipschitz dependence of $\alpha$ under dilation.}
Set $\lambda:=r_*^{-1}$. The symmetric difference $U_*\Delta\lambda U_*$ is
contained in an annulus of width $|\lambda-1|R$; by the BDV perimeter bound,
\begin{equation}\label{eq:dilL1}
|U_*\Delta\lambda U_*|\le R\,P(U_*)\,|\lambda-1|+o(|\lambda-1|)
\le 2RP_{\rm sel}(N,R)\,|\lambda-1|,
\end{equation}
for $|\lambda-1|\le1/2$. Combined with~\eqref{eq:42ii},
\[
|\alpha(\wt\Om)-\alpha(U_*)|\le C_2(R)\,|U_*\Delta\lambda U_*|
\le 2RC_2P_{\rm sel}\,|\lambda-1|.
\]
By Lemma~\ref{lem:2}, $|\lambda-1|\le 2C_5\delta$ when $|r_*-1|\le1/2$.
Hence
\[
|\alpha(\wt\Om)-\alpha(U_*)|\le C_\alpha(N,R)\,\delta,\qquad
C_\alpha:=4RC_2P_{\rm sel}C_5.
\]
Combined with Lemma~\ref{lem:1},
\begin{equation}\label{eq:alphatilde}
|\alpha(\wt\Om)-\eps|
\le|\alpha(\wt\Om)-\alpha(U_*)|+|\alpha(U_*)-\eps|
\le\eps\Bigl(C_\alpha q+\frac{\sqrt{2q+q^2}}{\tau}\Bigr)=\rho\eps.
\end{equation}
With $\tau=q^{1/4}$ and small $q$, $\rho\le2q^{1/4}+C_\alpha q\le1/2$.

\paragraph{Step 2: Energy scaling.}
For quasi-open $V\subset B_R$ and $\mu>0$ with $\mu V\subset B_R$,
\begin{equation}\label{eq:Escale}
E(\mu V)=\mu^{N+2}E(V).
\end{equation}
Indeed, if $u_V$ is the torsion function of $V$, then
$u_{\mu V}(x)=\mu^2 u_V(x/\mu)$ minimises
$\tfrac12\int|\nabla\cdot|^2-\int\cdot$ on $H^1_0(\mu V)$, and the change of
variable gives~\eqref{eq:Escale}.

\paragraph{Step 3: From $E(U_*)$ to $E(\wt\Om)$.}
Apply~\eqref{eq:Escale} with $\mu=\lambda=r_*^{-1}$:
\[
E(\wt\Om)-E(B_1)=\lambda^{N+2}\bigl(E(U_*)-E(B_1)\bigr)
+(\lambda^{N+2}-1)E(B_1).
\]
From Lemma~\ref{lem:1} and~\eqref{eq:Lem45},
\[
E(U_*)-E(B_1)\le\delta+\hat\eta^{-1}||U_*|-|B_1||\le(1+\hat\eta^{-1}C_4)\delta
=:C_6\delta.
\]
Since $|\lambda-1|\le1/2$, $\lambda^{N+2}\le(3/2)^{N+2}=:C_7(N)$ and
$|\lambda^{N+2}-1|\le C_8(N)C_5\delta$. Thus
\begin{equation}\label{eq:EtildeFinal}
E(\wt\Om)-E(B_1)\le(C_6C_7+C_8C_5|E(B_1)|)\,\delta=:C_E(N,R)\,\delta.
\end{equation}

\paragraph{Step 4: Quotient.}
By \eqref{eq:alphatilde} and \eqref{eq:EtildeFinal},
\[
Q_\alpha(\wt\Om)=\frac{E(\wt\Om)-E(B_1)}{\alpha(\wt\Om)}
\le\frac{C_E\delta}{(1-\rho)\eps}=\frac{C_E}{1-\rho}q
=\frac{C_E}{1-\rho}Q_\alpha(\Om_0).
\]
This proves~(d).

\section{The Selection Theorem}\label{sec:thm}

\begin{theorem}[Single-set quantitative selection]\label{thm:main}
Fix $N\ge2,R\ge2$. Let
$\alpha_{\rm sel}(N,R):=\eps_{\rm sel}(N,R)$,
$q_{\rm sel}(N,R):=\tau_{\rm reg}(N,R)^4$. Suppose $\Om_0\subset B_R$ is
quasi-open with $|\Om_0|=|B_1|$, $0<\eps:=\alpha(\Om_0)\le\alpha_{\rm sel}$,
$q:=Q_\alpha(\Om_0)\le q_{\rm sel}$. Set $\tau:=q^{1/4}$ and let $U_*$ be a
minimiser of $J$. Define $\wt\Om:=r_*^{-1}U_*$,
$r_*=(|U_*|/|B_1|)^{1/N}$. Then there exist $\rho=\rho(q)\le1/2$ and
$C_{\rm sel}(N,R)>0$ such that:
\begin{enumerate}
\item[\rm(i)] $|\wt\Om|=|B_1|$;
\item[\rm(ii)] $(1-\rho)\eps\le\alpha(\wt\Om)\le(1+\rho)\eps$, $\rho\le1/2$;
\item[\rm(iii)] $Q_\alpha(\wt\Om)\le\dfrac{C_{\rm sel}(N,R)}{1-\rho}\,q$;
\item[\rm(iv)] $\wt\Om$ inherits the BDV penalised-minimiser regularity
package (two-sided density estimates, Lipschitz/nondegeneracy of the
torsion potential, smooth Bernoulli coefficient);
\item[\rm(v)] the selected Bernoulli law on $\partial\wt\Om$ takes the
volume-rescaled form
\[
|\nabla\wt u(y)|^2=\wt\Lambda+\wt a_\sigma\,\wt\Psi_{r_*}(y),
\qquad\wt\Lambda=r_*^{-2}\Lambda,\quad\wt a_\sigma=r_*^{-1}a_\sigma.
\]
\end{enumerate}
\end{theorem}

\begin{proof}
Parts (i)--(iii) are Lemma~\ref{lem:3}. For (iv)--(v) see
Sections~\ref{sec:reg}--\ref{sec:bern}.
\end{proof}

\section{Free-Boundary Regularity Inherited from BDV}\label{sec:reg}

The minimiser $U_*=\{u_*>0\}$ satisfies the perturbed one-phase
quasi-minimality inequality (BDV (4.19) with $\sigma$ replaced by $\tau$):
for every $v\in H^1_0(B_R)$,
\begin{equation}\label{eq:qmin}
\tfrac12\!\int|\nabla u_*|^2-\!\int u_*+f_{\hat\eta}(|\{u_*>0\}|)
\le\tfrac12\!\int|\nabla v|^2-\!\int v+f_{\hat\eta}(|\{v>0\}|)
+C_R\,\tau\,\bigl|\{u_*>0\}\Delta\{v>0\}\bigr|,
\end{equation}
the error being BDV's bound for the variation of
$\sqrt{\eps^2+\tau^2(\alpha(U)-\eps)^2}$, dominated by
$\tau\,C_2(R)|U\Delta U_*|$.

For $\tau\le\tau_{\rm reg}(N,R)$ the BDV Alt--Caffarelli regularity package
applies with constants depending only on $N,R$ (Lemmas~4.9, 4.12, 4.16).
The volume-normalisation $\wt\Om=r_*^{-1}U_*$ is a bilipschitz dilation with
$|\lambda-1|\le1/2$; all density, Lipschitz, nondegeneracy and $C^k$
constants transfer with degradation at most $(3/2)^N$.

\section{The Volume-Normalised Bernoulli Law}\label{sec:bern}

BDV Lemma~4.16 gives on $\partial U_*$,
\begin{equation}\label{eq:Bern1}
|\nabla u_{U_*}(x)|^2=\Lambda+a_\sigma\,\Psi_{U_*}(x),
\qquad |a_\sigma|\le\tau,
\end{equation}
with $\Psi_{U_*}(x)=|x-x_{U_*}|-V_{U_*}\cdot x$ and
$V_{U_*}=\int_{U_*}\frac{z-x_{U_*}}{|z-x_{U_*}|}\,dz$.

\begin{lemma}[Rescaled Bernoulli law]\label{lem:bern}
Set $\wt u(y):=r_*^{-2}u_{U_*}(r_*y)$. Then $-\Delta\wt u=1$ in $\wt\Om$,
$\wt u=0$ on $\partial\wt\Om$, and on $\partial\wt\Om$
\begin{equation}\label{eq:Berntilde}
|\nabla\wt u(y)|^2=\wt\Lambda+\wt a_\sigma\,\wt\Psi_{r_*}(y),
\end{equation}
with $\wt\Lambda=r_*^{-2}\Lambda$, $\wt a_\sigma=r_*^{-1}a_\sigma$, and
\[
   \wt\Psi_{r_*}(y):=
   |y-r_*^{-1}x_{U_*}|-r_*^{N}V_{\wt\Om}\cdot y.
\]
This differs from the unscaled $\alpha$-bracket only by the harmless
factor $r_*^N=1+O(\delta)$ in the barycentric term.
\end{lemma}

\begin{proof}
$\nabla\wt u(y)=r_*^{-1}\nabla u_{U_*}(r_*y)$ and $\Delta\wt
u(y)=\Delta u_{U_*}(r_*y)=-1$. On $\partial\wt\Om$, $r_*y\in\partial U_*$,
so by~\eqref{eq:Bern1},
$|\nabla\wt u(y)|^2=r_*^{-2}\Lambda+r_*^{-2}a_\sigma\,\Psi_{U_*}(r_*y)$.
Change of variables in the centroid integral gives
$V_{U_*}=r_*^N V_{\wt\Om}$ and
$\Psi_{U_*}(r_*y)=r_*\bigl[|y-r_*^{-1}x_{U_*}|-r_*^N V_{\wt\Om}\cdot y\bigr]$.
Reorganising yields~\eqref{eq:Berntilde}. Since $|r_*-1|\le C_5\delta$ both
prefactors degrade by $1+O(\delta)$.
\end{proof}

\section{Summary of Constants}\label{sec:constants}

\begin{center}
\renewcommand{\arraystretch}{1.25}
\begin{tabular}{|l|l|l|}
\hline
Symbol & Source & Dependence \\\hline
$\hat\eta,C_4$ & BDV Lemma 4.5 & $N,R$ \\
$C_1$ & BDV Lemma 4.2(i) & $N$ \\
$C_2$ & BDV Lemma 4.2(ii) & $R$ \\
$\sigma_2,\sigma_3$ & BDV Lemmas 4.9, 4.16 & $N,R$ \\
$\tau_{\rm ex},\tau_{\rm reg}$ & this note & $N,R$ \\
$P_{\rm sel}$ & BDV Lemma 4.6 & $N,R$ \\
$C_5$ & Lemma~\ref{lem:2} & $N,R$ \\
$C_\alpha=4RC_2P_{\rm sel}C_5$ & Step 1 of Lemma~\ref{lem:3} & $N,R$ \\
$C_E$ & Step 3 of Lemma~\ref{lem:3} & $N,R$ \\
$C_{\rm sel}=C_E$ & final quotient & $N,R$ \\
$\eps_{\rm sel},q_{\rm sel}$ & smallness regime~\eqref{eq:Sml} & $N,R$ \\
$\rho(q)=2q^{1/4}+C_\alpha q$ & \eqref{eq:alphatilde} & $N,R$, $q$ \\
$\wt\Lambda=r_*^{-2}\Lambda,\ \wt a_\sigma=r_*^{-1}a_\sigma$ &
rescaled Bernoulli & $N,R$ \\\hline
\end{tabular}
\end{center}

All constants depend only on $(N,R)$ via named BDV constants. None is set
by a compactness or contradiction step.

\begin{thebibliography}{9}
\bibitem{BDV} L. Brasco, G. De Philippis, B. Velichkov,
\emph{Faber--Krahn inequalities in sharp quantitative form},
Duke Math. J. \textbf{164} (2015), 1777--1831 (arXiv:1306.0392).
\end{thebibliography}

\end{document}
